% ==============================================================================
% Appendix B — Supplementary Figures and Tables
% ==============================================================================

\chapter{Supplementary Material}
\section{Tables}

\begin{table}[htbp]
  \centering
  \caption{Indicative GNFR to IPCC (1996) sector codes}
  \label{tab:gnfr_ipcc_mapping}
  \small
  \setlength{\tabcolsep}{3pt}
  \begin{tabular}{ll}
    \hline
    IPCC code(s) & Description \\
    \hline
    \hline
    \multicolumn{2}{@{}c@{}}{\textit{Power Generation ; A\_PublicPower}} \vspace{0.1cm} \\
     1A1a & Energy industry  \\
     1A1 & Power generation  \\
    \hline
    \hline
    \multicolumn{2}{@{}c@{}}{\textit{Industrial processes ; B\_Industry}} \vspace{0.1cm} \\
     1A(1b,1c,5b1),1B(1b,2a5,2a6,2b5),2C1b & Oil refineries and transformation / industry combustion \\
     1A1b;1B2a5 & Petroleum refining and distribution \\
     1A2 & Combustion in manufacturing industry \\
     2A & Non-metallic mineral processes  \\
     2B & Chemical processes \\
     2\_3; 2B; 2D; 1A1b,1B2a5 & Industrial processes and product use\\
     2C(1a \ldots 1f,2) & Iron and steel\\
     2C3--2C5 & Aluminum, magnesium and steel production \\
     2G & Non-energy use of fuels \\
    \hline
    \hline
    \multicolumn{2}{@{}c@{}}{\textit{Residential ; C\_OtherStationary Comb}} \vspace{0.1cm}\\
    1A4 & Residential and other stationary combustion \\
    \hline
    \hline
    \multicolumn{2}{@{}c@{}}{\textit{Fugitive emissions from solid fuels ;D\_Fugitive}} \vspace{0.1cm}\\
    1B(1a,2a1 \ldots 2a4, 2c); 7A & Fuel exploitation \\
    --- & Manufacturing of solid fuels and fugitive emissions \\
    \hline
    \hline
    \multicolumn{2}{@{}c@{}}{\textit{Solvents ; E\_Solvents}}\vspace{0.1cm} \\
    3 & Solvents production and application \\
    \hline
    \hline
    \multicolumn{2}{@{}c@{}}{\textit{Road transportation ; F\_RoadTransport}}  \vspace{0.1cm}\\ 
    1A3b & Road transportation \\
    \hline
    \hline
    \multicolumn{2}{@{}c@{}}{\textit{Ships ; G\_Shipping}}\vspace{0.1cm} \\
    1A3d\_1C2 & Ships \\ 
    \hline
    \hline
    \multicolumn{2}{@{}c@{}}{\textit{Non-road ground transportation ; I\_OffRoad}} \vspace{0.1cm} \\
    1A3c\_1A3e & Non-road ground transportation \\
    \hline
    \hline
    \multicolumn{2}{@{}c@{}}{\textit{Solid Waste and Waste Water ; J\_Waste}} \vspace{0.1cm} \\
    6A, 6D & Waste incineration \\
    6B & Solid waste disposal \\
    6C & Waste water \\
    \hline
    \hline
    \multicolumn{2}{@{}c@{}}{\textit{Agriculture ; K\_AgriLivestock \& L\_AgriOther }}\vspace{0.1cm} \\  
     4B & Manure management (K)\\
     4A & Enteric fermentation (L) \\
     4C\_4D1\_4D2\_4D4 & Agricultural soils (L)  \\
     4F & Agricultural waste / field burning (L)\\
    
    
    \hline
  \end{tabular}
  \medskip

  \footnotesize
  See \cite{cams_glob_ant} and \cite{ipcc_efdb}.
\end{table}



\begin{table}[h]
\centering
\caption{CAMS-REG-ANT v8.1: total emissions by pollutant for GNFR K (AgriLivestock) and L (AgriOther), summing area and point sources. Units: kg yr$^{-1}$.}
\label{tab:cams_kl_totals}
\begin{tabular}{lrr}
\hline
Pollutant & K (AgriLivestock) & L (AgriOther) \\
\hline
CH$_4$ & $1.529 \times 10^{10}$ & $6.434 \times 10^{8}$ \\
CO & $0$ & $1.413 \times 10^{9}$ \\
NH$_3$ & $3.172 \times 10^{9}$ & $3.243 \times 10^{9}$ \\
NMVOC & $2.612 \times 10^{9}$ & $7.561 \times 10^{8}$ \\
NO$_x$ & $7.319 \times 10^{6}$ & $1.071 \times 10^{8}$ \\
PM$_{10}$ & $2.745 \times 10^{8}$ & $4.742 \times 10^{8}$ \\
PM$_{2.5}$ & $5.542 \times 10^{7}$ & $1.483 \times 10^{8}$ \\
SO$_x$ & $0$ & $6.351 \times 10^{6}$ \\
CO$_{2}$ (fossil) & $0$ & $1.099 \times 10^{10}$ \\
CO$_{2}$ (biofuel) & $0$ & $2.490 \times 10^{10}$ \\
\hline
\end{tabular}
\end{table}

\subsection{Tables for agriculutral sector}

\begin{table}[htbp]
  \centering
  \begin{threeparttable}
    \caption{Mineral nitrogen rates $R_k$ from \texttt{synthetic\_N\_rate.json} (key \texttt{rates}). Units: kg~N~ha$^{-1}$. }
    \label{tab:synthetic-n-rates}
    \scriptsize
    \setlength{\tabcolsep}{6pt}
    \begin{tabular}{@{}lrr@{\hspace{2em}}lr@{}}
      \toprule
      Category $k$ & $R_k$ & & Category $k$ & $R_k$ \\
      \midrule
      wheat & 116 & & rapeseed & 146 \\
      maize\_grain & 169 & & soya & 52 \\
      rice & 118 & & cotton & 101 \\
      barley & 99 & & tobacco & 29 \\
      rye & 29 & & other\_industrial & 144 \\
      oats & 43 & & dry\_pulse & 15 \\
      other\_cereal & 52 & & tomato & 125 \\
      potato & 77 & & other\_veg & 103 \\
      sugar\_beet & 123 & & fodder\_crop & 51 \\
      sunflower & 51 & & olive & 58 \\
      vineyard & 28 & & fruit\_tree & 67 \\
      citrus & 67 & & nuts\_tree & 56 \\
      temporary\_grass & 90 & & permanent\_grass & 43 \\
      shrub\_grass & 10 & & fallow & 0 \\
      \bottomrule
    \end{tabular}
    
    \begin{tablenotes}
      \footnotesize
      \item \textit{Note:} This data was derived from \citet{einarsson_2021}. The data can be accessed \href{https://figshare.com/articles/dataset/EuropeAgriDB_v1_0_Source_code_and_data/16540314}{here} and the code \texttt{synth\_rate\_n.py} creates the tables presented above.
    \end{tablenotes}
  \end{threeparttable}
\end{table}

\begin{table}[htbp]
\centering
\begin{threeparttable}
\caption{NH$_3$ housing emission factors used for livestock categories. Values in gN head$^{-1}$ day$^{-1}$ are derived from \citet{Misselbrook_2000} using standard liveweight assumptions consistent with IPCC guidelines \cite{IPPC_2006_V10}.}
\label{tab:livestock_housing_factors}
\begin{tabular}{lcccc}
\toprule
Animal & gN/LU/day & Liveweight (kg) & gN/head/day & IPCC table \\
\midrule
Bovine  & 34.3  & 600  & 41.16 & T10A-4 \\
Pigs    & 80.0  & 50   & 8.00  & T10A-7 \\
Poultry & 150.0 & 1.8  & 0.54  & T10A-9 \\
Sheep   & 3.0   & 48.5 & 0.291 & T10A-9 \\
Goats   & 3.0   & 38.5 & 0.231 & T10A-9 \\
\bottomrule
\end{tabular}
\begin{tablenotes}[flushleft]
\footnotesize
\item Per-head emission factors are obtained as
\[
EF_{\mathrm{animal}} \; [\mathrm{gN \, head}^{-1}\mathrm{day}^{-1}]
=
EF_{\mathrm{animal}} \; [\mathrm{gN \, LU}^{-1}\mathrm{day}^{-1}]
\times \frac{\mathrm{Liveweight}}{500},
\]
assuming the conventional definition \(1\ \mathrm{LU} = 500\ \mathrm{kg}\) liveweight.
\end{tablenotes}
\end{threeparttable}
\end{table}

\begin{table}[htbp]
  \centering
  \begin{threeparttable}
  \caption{Eligible crop types for the biomass-burning proxy (\(\mathcal{C}_{\mathrm{emep}}\), NFR~4.F).}
  \label{tab:emep-s-values}
  \small
  \begin{tabular}{@{}llcl@{}}
    \hline
    Crop & LC1 & \(s\) & Note \\
         &      & (kg/kg) &      \\
    \hline
    Wheat (common)        & B11 & 1.3 & \\
    Wheat (durum)         & B12 & 1.3 & a \\
    Barley                & B13 & 1.2 & \\
    Rye                   & B14 & 1.6 & \\
    Oats                  & B15 & 1.3 & \\
    Maize (grain)         & B16 & 1.0 & \\
    Rice                  & B17 & 1.4 & \\
    Triticale             & B18 & 1.2 & b \\
    Other cereals         & B19 & 1.2 & c \\
    Soya                  & B33 & 2.1 & \\
    Dry pulses            & B41 & 1.7 & d \\
    \hline
    \multicolumn{4}{@{}l@{}}{\textit{Non-eligible crops (\(s = 0\))}} \\
    Root crops            & B21--B23          & 0 & \\
    Oilseeds excl.\ soya  & B31, B32, B34--B37 & 0 & \\
    Vegetables, flowers   & B42--B45          & 0 & \\
    Fodder crops          & B51--B55          & 0 & \\
    Permanent crops       & B71--B84          & 0 & \\
    Fallow, grassland     & BX1, BX2, E       & 0 & \\
    \hline
  \end{tabular}
  \begin{tablenotes}[flushleft]
    \footnotesize
    \item Source: residue-to-yield ratios from EMEP/EEA Guidebook 2023, Table~3-2; LUCAS 2022 LC1 codes from \texttt{Agriculture/config/synthetic\_N\_rate.json}. All LC1 codes not listed as eligible receive \(s = 0\).
    \item[a] Same \(s\) as common wheat.
    \item[b] Not reported in EMEP Table~3-2; assigned the barley value.
    \item[c] EU production-weighted cereal mean.
    \item[d] Mean of peas (\(1.5\)) and beans (\(2.1\)).
  \end{tablenotes}
  \end{threeparttable}
\end{table}

\begin{table}[htbp]
  \centering
  \begin{threeparttable}
  \caption{NMVOC-emitting crop categories and assigned per-hectare emission factors.}
  \label{tab:nmvoc-ef}
  \small
  \begin{tabular}{@{}p{4.8cm}p{2.4cm}ccp{3.2cm}@{}}
    \toprule
    Crop category (Th\"unen) & LUCAS LC1 & EF class & EF & Note \\
                             &           &          & (kg ha$^{-1}$ yr$^{-1}$) & \\
    \midrule
    Wheat                    & B11, B12  & wheat & 0.32 & a \\
    Rye                      & B14       & rye   & 1.03 & a \\
    Oilseed rape             & B32       & rape  & 1.34 & b \\
    Clover, grass mixtures, alfalfa & B51--B55 & grass & 0.41 & c \\
    Pastures and meadows     & E10, E20  & grass & 0.41 & c \\
    \addlinespace
    Barley                   & B13       & wheat & 0.32 & d \\
    Oat                      & B15       & wheat & 0.32 & d \\
    Triticale                & B18       & wheat & 0.32 & d \\
    Cereals, whole-plant harvest & B19   & wheat & 0.32 & d \\
    Grain maize              & B16       & wheat & 0.32 & d \\
    Maize for silage         & B16       & wheat & 0.32 & d \\
    Rice                     & B17       & wheat & 0.32 & d \\
    Tuber crops              & B21--B23  & wheat & 0.32 & d \\
    Pulses                   & B41       & wheat & 0.32 & d \\
    \midrule
    \multicolumn{5}{@{}l@{}}{\textit{Non-emitting categories (\(EF = 0\))}} \\
    Permanent crops          & B71--B84  & --    & 0    & e \\
    Other oilseeds           & B31, B33--B37 & -- & 0   & e \\
    Vegetables and flowers   & B42--B45  & --    & 0    & e \\
    Fallow                   & BX1, BX2  & --    & 0    & e \\
    Shrub grassland          & E30       & --    & 0    & e \\
    \bottomrule
  \end{tabular}
  \begin{tablenotes}[flushleft]
    \footnotesize
    \item Source: three base EFs from EMEP/EEA Guidebook Table~3.3 \citep{EEA_guidebook}; crop assignment follows the Th{\"u}nen NIR methodology (Section~2.8, Table~11 in \citet{thuenen_nmvoc}).
    \item[a] Measured EF reported for the crop class in the Guidebook (K\"onig/Lamb).
    \item[b] Measured EF reported for rape (K\"onig, 1995).
    \item[c] Grass EF at 15\textdegree C.
    \item[d] No crop-specific EF available; wheat EF assigned following the Th{\"u}nen categorization.
    \item[e] Categories treated as non-emitting in this proxy.
  \end{tablenotes}
  \end{threeparttable}
\end{table}


\section{GNFR~C residential and commercial downscaling (reference tables)}
\label{app:rescom_downscaling}

This section complements Section~\ref{sec:rescom}. Eurostat row strings follow
the custom extract labelling; sheet names may differ slightly between downloads
and are set in the run configuration.

\begin{table}[htbp]
  \centering
  \caption{Aggregated model classes $k$ for GAINS-to-proxy mapping (GNFR~C).}
  \label{tab:rescom_model_classes}
  \begin{tabular}{p{3.5cm}p{9.5cm}}
    \hline
    Class & Intended use \\
    \hline
    \texttt{R\_FIREPLACE} & Residential fireplaces (household or agriculture
      forestry other sector in GAINS) \\
    \texttt{R\_HEATING\_STOVE} & Residential heating stoves \\
    \texttt{R\_COOKING\_STOVE} & Cooking stoves, three-stone stoves \\
    \texttt{R\_BOILER\_MAN} & Single-house or medium boilers, manual firing,
      non-commercial contexts \\
    \texttt{R\_BOILER\_AUT} & Single-house or medium boilers, automatic firing,
      non-commercial contexts \\
    \texttt{C\_BOILER\_MAN} & Medium boilers, manual, commercial \\
    \texttt{C\_BOILER\_AUT} & Medium boilers, automatic, commercial \\
    \hline
  \end{tabular}
\end{table}

\begin{table}[htbp]
  \centering
  \caption{Eurostat household end-use row labels (Table~3) used to match
    national shares; short synonyms may also be configured.}
  \label{tab:rescom_eurostat_table3}
  \footnotesize
  \begin{tabular}{@{}p{0.92\textwidth}@{}}
    \hline
    Label (full string) \\
    \hline
    Final consumption - other sectors - households - energy use - space heating \\
    Final consumption - other sectors - households - energy use - space cooling \\
    Final consumption - other sectors - households - energy use - water heating \\
    Final consumption - other sectors - households - energy use - cooking \\
    Final consumption - other sectors - households - energy use - lighting and electrical appliances \\
    Final consumption - other sectors - households - energy use - other end use \\
    \hline
  \end{tabular}
\end{table}

\begin{table}[htbp]
  \centering
  \caption{Eurostat Table~1 fuel column headers (custom extract) for
    residential fuel mix; reserved for future weighting of national structure.}
  \label{tab:rescom_eurostat_table1}
  \footnotesize
  \begin{tabular}{@{}p{0.92\textwidth}@{}}
    \hline
    Column header \\
    \hline
    Solid fossil fuels, peat, peat products, oil shale and oil sands \\
    Natural gas \\
    Oil and petroleum products \\
    Renewables and biofuels \\
    Electricity \\
    Heat \\
    \hline
  \end{tabular}
\end{table}

\begin{table}[htbp]
  \centering
  \caption{Illustrative GAINS fuel $\rightarrow$ EMEP \texttt{tables.fuel}
    substring hints used to break ties in emission-factor row matching.}
  \label{tab:rescom_emep_fuel_hints}
  \footnotesize
  \begin{tabular}{p{4.5cm}p{7.5cm}}
    \hline
    GAINS fuel contains & EMEP fuel field contains \\
    \hline
    Natural gas & Natural gas \\
    Gas oil & Gas oil \\
    LPG & Natural gas \\
    Brown coal / Hard coal / Derived coal & Coal fuels \\
    lignite & Coal fuels \\
    Biogas & Natural gas \\
    \hline
  \end{tabular}

  \medskip
  \footnotesize
  Additional hints may be added for biomass-rich fuels once guidebook rows are
  aligned explicitly; ambiguous solid-fuel rows should be checked against
  \texttt{EMEP\_emission\_factors.json}.
\end{table}

\section{GNFR~E solvent downscaling (reference tables)}
\label{app:solvents_downscaling}

This section complements Section~\ref{sec:solvents}. It summarises the
archetype construction and the subsector-to-archetype mixing adopted for
the solvent and product use block.

\begin{table}[htbp]
  \centering
  \caption{Archetype activity surfaces and component weights used for
    GNFR~E solvent downscaling.}
  \label{tab:solvents_archetypes}
  \footnotesize
  \begin{tabular}{p{2.2cm}p{8.8cm}}
    \hline
    Archetype & Components and weights \\
    \hline
    household &
    population (0.70), residential urban share (0.30) \\
    service &
    urban fabric (0.40), service-oriented land (0.35), service land use from
    OpenStreetMap (0.25) \\
    industrial &
    industrial land from CORINE (0.35), industrial land use from
    OpenStreetMap (0.35), industrial building proxy (0.30) \\
    infrastructure &
    road length (0.35), weighted road length by road hierarchy (0.35),
    transport-related area (0.20), roof or built-surface proxy (0.10) \\
    \hline
  \end{tabular}
\end{table}

\begin{table}[htbp]
  \centering
  \caption{Nine solvent subsectors and their $\beta$ weights over the
    four archetype activity surfaces.}
  \label{tab:solvents_subsectors_beta}
  \footnotesize
  \begin{tabular}{p{3.2cm}cccc}
    \hline
    Subsector & household & service & industrial & infrastructure \\
    \hline
    2D3a domestic solvent use & 1.00 & 0.00 & 0.00 & 0.00 \\
    2D3b road paving & 0.00 & 0.00 & 0.00 & 1.00 \\
    2D3c asphalt roofing & 0.00 & 0.00 & 0.20 & 0.80 \\
    2D3d coating application & 0.40 & 0.20 & 0.40 & 0.00 \\
    2D3e degreasing and cleaning & 0.00 & 0.30 & 0.70 & 0.00 \\
    2D3f dry cleaning & 0.00 & 1.00 & 0.00 & 0.00 \\
    2D3g chemical products & 0.00 & 0.00 & 1.00 & 0.00 \\
    2D3h printing & 0.00 & 0.70 & 0.30 & 0.00 \\
    2D3i other solvent use & 0.50 & 0.30 & 0.20 & 0.00 \\
    \hline
  \end{tabular}
\end{table}

\section{GNFR~J waste downscaling (reference tables)}
\label{app:waste_downscaling}

This section complements Section~\ref{sec:waste}. It summarises the
operational family aggregation adopted for GNFR~J and the main spatial
datasets used to build the pollutant-specific waste weighting rasters.

\begin{table}[htbp]
  \centering
  \caption{Operational proxy families used for GNFR~J waste downscaling
    and the underlying CEIP/NRD subsectors.}
  \label{tab:waste_families}
  \footnotesize
  \begin{tabular}{p{3.4cm}p{4.1cm}p{5.0cm}}
    \hline
    Proxy family & CEIP/NRD subsectors & Interpretation \\
    \hline
    Solid waste treatment &
    5A, 5B1, 5B2 &
    Solid waste disposal on land, composting, and anaerobic digestion on biogas facilities \\
    Wastewater &
    5D1, 5D2, 5D3 &
    Domestic, industrial, and other wastewater handling \\
    Residual diffuse waste &
    5C2, 5E &
    Open waste burning and residual ``other waste'' processes \\
    \hline
  \end{tabular}

  \medskip
  \footnotesize
  Incineration-like subsectors 5C1a and 5C1b(i--vi) are excluded from the
  default area-weighting workflow to avoid double counting with point-source
  treatment of waste incineration facilities.
\end{table}

\begin{table}[htbp]
  \centering
  \caption{Spatial proxy families and principal geospatial support layers
    used for GNFR~J waste weighting.}
  \label{tab:waste_proxy_inputs}
  \footnotesize
  \begin{tabular}{p{3.4cm}p{4.5cm}p{4.6cm}}
    \hline
    Proxy family & Main spatial inputs & Intended role \\
    \hline
    Solid waste treatment &
    CORINE class~132, CORINE class~121, optional OpenStreetMap waste features &
    Landfill-type support with secondary industrial/commercial support and optional geometric refinement \\
    Wastewater &
    UWWTD agglomerations, UWWTD treatment plants, GHSL population, Copernicus imperviousness, optional industrial land &
    Settlement- and infrastructure-based support for domestic and industrial wastewater handling \\
    Residual diffuse waste &
    GHSL population, GHSL SMOD / rural settlement classes, Copernicus imperviousness &
    Diffuse settlement-related support with added emphasis on rural and low-density population \\
    \hline
  \end{tabular}
\end{table}

\section{Additional Emission Maps}
\label{app:emission_maps}

\todo[inline]{Place supplementary emission maps here (e.g., all GNFR sectors,
both cities, all pollutants) that are too detailed for the main body.}

\section{Model Evaluation Statistics}
\label{app:eval_stats}

\todo[inline]{Include extended evaluation tables for all monitoring stations and
pollutants evaluated.}

\section{Scenario Emission Budgets}
\label{app:scenario_budgets}

\todo[inline]{Provide full emission budget tables for each scenario compared to
the baseline, for both cities.}
